\scanpage{176}\mspage{31}%
haut s'interprète immédiatement comme
\[
\varphi(F_1)\simeq p_*{\unsure{J|j}}^*(\widetilde F_1)
\]
où $F_1\mapsto\widetilde F_1$ est la restriction de $X_1$ à $\widetilde X_1$ (qui
est donc une équivalence de catégories pour les
faisceaux (loc. const.) envisagés). \deleted{\unsure{Il|Il s'|\ensuremath{J}} \unsure{peut|peut se|\ensuremath{J^*}}} \inserted{La donnée}
d'un\inserted{\unsure{e}} $F=(F_0,F_1,u:F_0\to\varphi F_1)$ peut donc aussi
s'expliciter comme la donnée de
\[
F_0,\ \widetilde F_1,\ \widetilde u:p^*(F_0)\longrightarrow J^*(\widetilde F_1)
\]
où $F_0$ ($\widetilde F_1$) est un faisceau loc. constant sur
$X_0$ (resp. $\widetilde X_1$). Il faut se rappeler de plus
qu'ici $p$ est une vraie fibration, et $J$ une inclusion
(et pratiquement, -- cas $X_0,X_1$ lisses --
\deleted{\unsure{une}} l'inclusion du bord dans une variété
à bord).

Si on préfère, on peut aussi prendre le
diagramme un peu moins joli (mais un
peu plus canonique)
\diagram{\begin{tikzcd}[ampersand replacement=\&,column sep=large,row sep=large]
\mathring T \arrow[r,"J'"] \arrow[d,"p'"] \& X_1\\
% Small isolated bent mark at whole-image x=350--357, y=1841--1850,
% about 380 px left of X_0; uncertain origin, no textual or mathematical role assigned.
\raisebox{0pt}[0pt][0pt]{\parbox{0pt}{\figure{\begin{tikzpicture}[overlay]
\draw[line width=.3pt] (-4.8cm,.04cm) .. controls (-4.77cm,.01cm) and (-4.8cm,-.01cm) .. (-4.77cm,-.04cm);
\end{tikzpicture}}}}X_0 \&
\end{tikzcd}}
revenant essentiellement au $\overline m$, car formé d'espaces
homotopes au précédent. On peut $\overline m$ y remplacer
$X_0$ par $T$ ($X_0$ en étant un rétract par déformation) et écrire
\diagram{\begin{tikzcd}[ampersand replacement=\&,column sep=large,row sep=large]
\mathring T \arrow[r,hook,"J'"] \arrow[d,hook,"p''"] \& X_1\\
\unsure{T|I} \&
\end{tikzcd}}
% archive sheet number: 176
% The bottom target is cut close to the paper edge; no content supplied from another page.
